\documentclass[a4paper,slidestop,xcolor=pst,dvips,blue]{beamer}

\usepackage{beamerthemesplit}
\usepackage[utf8]{inputenc}
\usepackage[spanish]{babel}
\usepackage{graphicx}
\usepackage{pstricks} % PSTricks package
\usepackage{setspace}
\usepackage{multirow}
\usepackage{listings}
\usepackage{pgfpages}
\usepackage{hyperref}
\usepackage{etoolbox}
\usepackage{epstopdf}

\makeatletter
\patchcmd{\beamer@sectionintoc}{\vskip1.5em}{\vskip0.5em}{}{}
\makeatother

\setbeamercovered{dynamic}
\setcounter{tocdepth}{2}
\setbeamercolor{frametitle}{fg=black,bg=white}
\setbeamercolor{section in toc shaded}{fg=black}
\setbeamercolor{section in toc}{fg=red}
\setbeamercolor{subsection in toc shaded}{fg=black}
\setbeamercolor{subsection in toc}{fg=red}
\setbeamerfont{section in toc}{size=\small}
\setbeamerfont{subsection in toc}{size=\small}
\setbeamertemplate{section in toc shaded}[default][99]
\setbeamertemplate{subsection in toc shaded}[default][99]

\AtBeginSection[]
{\begin{frame}[c]
  \frametitle{Índice}
	\tableofcontents[currentsection,
        sectionstyle=show/shaded,
        subsectionstyle=hide]
\end{frame}}

\AtBeginSubsection[]
{\begin{frame}[c]
	\frametitle{Índice}
	\tableofcontents[
  		currentsection,
  		sectionstyle=shaded/shaded,
  		currentsubsection,
  		subsectionstyle=show/shaded/hide
		]
\end{frame}}

\setbeamercolor{frametitle}{fg=black,bg=white}

\setbeamertemplate{frametitle}{
	\begin{centering}
		\insertframetitle
		\par
	\end{centering}
}

\usetheme[secheader]{Boadilla} 

\title[Ing. Requisitos]{Ingeniería de Requisitos}

\author[Pablo Sánchez]{\alert{Pablo Sánchez}}

\institute[I2E]{
		   Dpto. Ingeniería Informática y Electrónica \\
		   Universidad de Cantabria \\
		   Santander (Cantabria, España) \\
		   p.sanchez@unican.es
}

\date{}

\begin{document}

\begin{frame}[c]
	\titlepage
	\begin{columns}
		\column{0.50\linewidth}
			\centering
			\includegraphics[width=.28\textwidth,keepaspectratio=true]{images/istr.eps}
		\column{0.50\linewidth}
			\centering
			\includegraphics[width=.25\textwidth,keepaspectratio=true]{images/uc.eps}
    \end{columns}
\end{frame}

\begin{frame}[c]
	\frametitle{Índice}
	\begin{enumerate}
        \item \alert{Índice}
        \item Datos Administrativos
        \item Objetivos y Temario
        \item Metodología
        \item Evaluación
        \item Bibliografía
	\end{enumerate}
\end{frame}

\section{Datos Administrativos}

\subsection{Profesorado}

\begin{frame}[c]
	\frametitle{Profesorado}
	\begin{center}
		\alert{Pablo S\'{a}nchez Barreiro}  \\
		\begin{small}
		Despacho 1069 \\
		Departamento de Ingeniería Informática y Electrónica \\
		p.sanchez@unican.es \\
		\end{small}
	\end{center}
\end{frame}

\subsection{Horarios}

\begin{frame}[c]
	\frametitle{Horario Clases}
	\begin{small}
	\begin{center}
	\begin{tabular}{||l|c|c|c|c|c||}
	\hline \hline
				   & Lunes  & Martes      & Miércoles    & Jueves       & Viernes      \\ \hline \hline
    08:30 - 09:30  &        &             &              &              &              \\ \hline
    09:30 - 10:30  &        &             &              &              & Laboratorio  \\ \hline \hline
	10:45 - 11:45  &        &             &              &              & Laboratorio  \\ \hline
    11:45 - 12:45  &        &  Aula       & Aula         &              &              \\ \hline
    12:45 - 13:45  &        &             &              &              &              \\ \hline
	\end{tabular}
	\end{center}
	\end{small}
	\begin{itemize}
		\item<2-> La asistencia a las clases teóricas y prácticas no es obligatoria, \uncover<3->{\alert{pero si altamente recomendable e incluso necesaria}.}
	\end{itemize}
\end{frame}

\begin{frame}[c]
	\frametitle{Horario Tutorías}
    \begin{itemize}[<+->]
        \item A cualquier hora, preferentemente de mañana.
        \item Disponible tardes de Lunes y Jueves.
        \item Si se quiere asegurar disponibilidad, avisar con antelación.
        \item Disponible a través de correo y \alert{moodle}.
	\end{itemize}
\end{frame}

\section{Objetivos y Temario}

\subsection{Objetivos}

\begin{frame}[t]
	\frametitle{Objetivos de la Asignatura}
    \begin{enumerate}
			\item<1-> Conocer y comprender la importancia y papel de la \emph{Ingeniería de Requisitos}.
            \item<2-> Conocer, comprender y saber diseñar y ejecutar procesos de \emph{Ingeniería de Requisitos}.
			\item<3-> Conocer, comprender y saber aplicar técnicas para la \emph{Captura de Requisitos}.
			\item<4-> Conocer, comprender y saber aplicar técnicas para la especificación de \emph{Requisitos Funcionales}.
			\item<5-> Conocer, comprender y saber aplicar técnicas para la especificación de \emph{Requisitos No Funcionales}.
			\item<6-> Conocer, comprender y saber aplicar técnicas para la especificación de \emph{Procesos de Negocio}.
	 \end{enumerate}
\end{frame}

\subsection{Resultados de Aprendizaje}

\begin{frame}[t]
	\frametitle{Resultados de Aprendizaje}
    	\begin{enumerate}
			\item<1-> El alumno conocerá, comprenderá y sabrá utilizar la terminología propia de la \emph{Ingeniería de Requisitos}.
			\item<2-> El alumno sabrá diseñar y ejecutar procesos de \emph{Ingeniería de Requisitos}.
            \item<3-> El alumno sabrá diseñar y ejecutar procesos de \emph{Captura de Requisitos}, empleando diferentes técnicas y estrategias, y sabiendo escoger la más adecuada en función de cada tipo de proyecto.
            \item<4-> El alumno sabrá especificar requisitos funcionales utilizando para ello \emph{objetivos}, \emph{escenarios}, \emph{casos de uso} e \emph{historias de usuario}.
            \item<5-> El alumno sabrá especificar requisitos no funcionales, en especial los requisitos relacionados con la seguridad, utilizando para ello \emph{objetivos} y catálogos predefinidos de requisitos no funcionales.
            \item<6-> El alumno sabrá modelar procesos de negocio mediante alguna notación estándar, como BPMN.
		\end{enumerate}
\end{frame}

\subsection{Temario}

\begin{frame}[c]
	\frametitle{Temario}
	\begin{enumerate}[<+->]
		\item Procesos de Ingeniería de Requisitos.
		\item Técnicas de Captura de Requisitos.
		\item Especificación y Modelado de Requisitos Funcionales.
		\item Especificación y Modelado de Requisitos No Funcionales.
		\item Especificación y Modelado de Procesos de Negocio.
	\end{enumerate}
\end{frame}

\section{Metodología}

\subsection{Plataforma}

\begin{frame}[c]
	\frametitle{Plataforma de Trabajo}
	\begin{itemize}
		\item<1-> La plataforma de trabajo de la asignatura es \emph{moodle}.
		\item<2-> Todas las notificaciones y publicaciones se harán a través de \emph{moodle}.
		\item<3-> \alert{Es obligación del alumno estar atento a las posibles notificaciones y avisos que se realicen a través de moodle}.
	\end{itemize}
\end{frame}

\subsection{Actividades}

\begin{frame}
	\frametitle{Clases en Aula}
	\begin{block}{Objetivo}
        Entender los conocimientos teóricos que constituyen la base de las habilidades, destrezas y competencias a adquirir al final de la asignatura.
	\end{block}
	\begin{itemize}
        \item<2-> Preparación autónoma y exposición del contenido teórico.
        \item<3-> Actividades de aprendizaje activo en función del contenido y predisposición del alumnado.
		\item<4-> Resolución autónoma, dirigida y supervisada de ejercicios.
		\item<5-> Clases magistrales puntuales usando pizarra y/o transparencias.
		\item<6-> NOTA: \alert{Las transparencias} son material de apoyo a la docencia, \alert{no son apuntes}.
		\item<7-> \alert{Sin conocimiento teórico es imposible alcanzar las habilidades prácticas}.
	\end{itemize}
\end{frame}

\begin{frame}[c]
	\frametitle{Clases en Laboratorio}
	\begin{block}{Objetivo}
		Aplicar los conceptos teóricos aprendidos en las clases de aula al desarrollo de un sistema software real de mediana escala, con el objetivo
        de desarrollar las competencias procedimentales y actitudinales deseadas.
	\end{block}
	\begin{itemize}
        \item<2-> Desarrollo de un proyecto de especificación de requisitos, dividido en varias etapas.
		\item<3-> Dos etapas diferenciadas: \emph{definición del proceso} y \emph{ejecución del proceso} de Ingeniería de Requisitos.
        \item<4-> La primera parte del proyecto requiere del trabajo en equipo y de actividades de grupo.
        \item<5-> La segunda parte se realizará de manera coordinada pero individual.
	\end{itemize}
\end{frame}

\section{Métodos de Evaluación y Calificación}

\subsection{Fórmula Calificación}

\begin{frame}[c]
	\frametitle{Cálculo de la Calificación Final}
	\begin{block}{Fórmula de Cálculo de la Calificación Final}
		\begin{scriptsize}
        \begin{tabular}{ll}
			Calificación Final = & SI Calificación Proyecto >= 5.00 ENTONCES \\
                                 & \ \ \ SI Calificación Prueba Final  $>$= Calificación Proyecto ENTONCES \\
                                 & \ \ \ \ \ \ Calificación Prueba Final + 1  \\
                                 & \ \ \ SINO \\
                                 & \ \ \ \ \ \ Calificación Proyecto \\
                                 & \ \ \ FINSI \\
                                 & SINO \\
                                 & \ \ \  Calificación Prueba Final \\
                                 & FINSI \\
                                 & + Calificación Participación
		\end{tabular}
        \end{scriptsize}
	\end{block}
	\begin{itemize}
		\item<2-> Es necesario presentarse a la prueba final para superar la asignatura.
        \item<4-> La calificación por participación es un extra adicional que variará entre (0, 0.5).
        \item<5-> La calificación por participación dependerá del desempeño y la actitud de cada alumno en el aula.
	\end{itemize}
\end{frame}

\subsection{Métodos y Criterios de Evaluación}

\begin{frame}
	\frametitle{Prueba Evaluable Escrita Final}
	\begin{itemize}[<+->]
		\item Prueba evaluable escrita en las fechas asignadas por el centro.
		\item La prueba contendrá ejercicios y cuestiones sobre razonamientos teóricos.
		\item Tendrá una duración de unas 4 horas, dividida en dos partes.
        \item Se podrá hacer uso todo tipo de material escrito.
        \item En ningún caso se dejará hacer uso de dispositivos electrónicos, especialmente, de aquellos con capacidades de comunicación inalámbrica.
		\item \alert{El material escrito debe servir para consultar cuestiones puntuales, pero en el caso ideal no debería hacerse ningín uso de los mismos}.
		\item Aquellas respuestas que sean una simple copia de lo disponible en algún tipo de material de la asignatura se calificarán con 0 puntos.
        \item Hay modelos de pruebas escritas disponibles en moodle.
	\end{itemize}
\end{frame}

\begin{frame}[c]
	\frametitle{Evaluación del Proyecto}
	\begin{itemize}[<+->]
        \item La calificación de la primera parte es única para todo el grupo.
        \item La calificación de la segunda parte es individual.
        \item Calificación final determinada por la media ponderada de las diversas entregas.
        \item Por cada entrega se proporcionarán unos criterios de evaluación y calificación claros.
        \item Hay que demostrar un mínimo de desempeño en la parte individual para poder superar la asignatura.
	\end{itemize}
\end{frame}

\begin{frame}[c]
	\frametitle{Peso de las Etapas del Proyecto}
    \begin{center}
	\begin{tabular}{||l|r||}
    \hline \hline
    Actas y Planes de Trabajo        & 5\% \\ \hline
    Definición de la Visión del Proyecto        & 10\% \\ \hline
    Definición del Contexto                           & 15\% \\ \hline
    Identificación de las Fuentes de Requisitos & 10\% \\ \hline
    Diseño del Plan de Captura de Requisitos    & 10\% \\ \hline \hline
    Diseño de una Actividad de Captura de Requisitos & 10\% \\ \hline
    Especificación de un Modelo de Objetivos    & 15\% \\ \hline
    Especificación de Requisitos Funcionales    & 15\% \\ \hline
    Especificación de Requisitos No Funcionales & 10\% \\ \hline \hline
    Modelado de Procesos de Negocio             &  5\% \\ \hline \hline
	\end{tabular}
\end{center}
\end{frame}

\section{Bibliografía}

\begin{frame}[c]
	\frametitle{Bibliografía}
\begin{thebibliography}{}

\bibitem[\protect\astroncite{Pohl}{2010}]{pohl:2010}
Pohl, K. (2010).
\newblock {\em {Requirements Engineering: Fundamentals, Principles and
  Techniques}}.
\newblock Springer.

\bibitem[\protect\astroncite{Gray et~al.}{2012}]{gray:2012}
Gray, D., Brown, S., and Macanufo, J. (2012).
\newblock {\em {Gamestorming: 83 juegos para innovadores, inconformistas y
  generadores del cambio}}.
\newblock Deusto.

\bibitem[\protect\astroncite{van Lamsweerde}{2009}]{lamsweerde:2009}
van Lamsweerde, A. (2009).
\newblock {\em {Requirements Engineering: From System Goals to UML Models to
  Software Specifications}}.
\newblock Wiley.

\end{thebibliography}
\end{frame}

\begin{frame}[c]
	\frametitle{Bibliografía}
\begin{thebibliography}{}

\bibitem[\protect\astroncite{Cockburn}{2000}]{cockburn:2000}
Cockburn, A. (2000).
\newblock {\em {Writing Effective Use Cases}}.
\newblock Addison-Wesley.

\bibitem[\protect\astroncite{Cohn}{2004}]{Cohn2004}
Cohn, M. (2004).
\newblock {\em {User Stories Applied: For Agile Software Development}},
\newblock Addison-Wesley Professional.

\bibitem[\protect\astroncite{Chung et~al.}{1999}]{chung:1999}
Chung, L., Nixon, B.~A., Yu, E., and Mylopoulos, J. (1999).
\newblock {\em {Non-Functional Requirements in Software Engineering}}.
\newblock Kluwer Academic Puvlishers.

\end{thebibliography}

\end{frame}

\end{document}
